\documentclass[11pt]{article}
\usepackage[a4paper,margin=25mm]{geometry}
\usepackage{amsmath,amssymb,amsthm,hyperref}
\newtheorem{lemma}{Lemma}\newtheorem{theorem}{Theorem}
\newcommand{\PP}{\mathcal P}
\title{Uniform dependence on the number of controlled prefix moments}
\author{Complete proof; three independent mathematical reviews completed}
\date{30 September 2026}
\begin{document}\maketitle
This note makes the fixed-parameter constants in
\path{../rational_designs/asymptotic/multiple_prefix_moment_conditioning.tex}
explicit enough for a growing-number-of-dice audit. It does not alone
supply a uniform complexity bound for the entire construction.
All norms are Lebesgue $L^2[0,1]$ norms.

\begin{lemma}[A uniform multiplier bound]\label{multiplier}
Let $Q$ be a real polynomial of degree $h$, with all roots in $[0,1]$
and leading coefficient of absolute value at least one. Then for
$p\in\PP_r$,
\[
 \|Qp\|_2\ge 2^{-1/2}[4h(r+1)^2]^{-h}\|p\|_2
\]
when $h\ge1$.
\end{lemma}
\begin{proof}
Remove intervals of radius $\varepsilon=[4h(r+1)^2]^{-1}$ around
the roots, counted with multiplicity. Their total length is at most
$2h\varepsilon$. The polynomial evaluation inequality
$\|p\|_\infty\le(r+1)\|p\|_2$ shows that at most half of its squared
norm is lost. On the remainder, factorization gives
$|Q|\ge\varepsilon^h$. Integrating there proves the claim.
No root-separation bound is needed.
\end{proof}

\begin{theorem}\label{conditioning}
There is an absolute constant $C$ with the following property.
Use $q\ge1$ Legendre bands and cumulative moments $M_0,\ldots,M_q$
from the cited multiple-prefix note, with multiplier degree $r$, total
band degree $D$, and $0<\eta\le1/(2q)$. Then
\[
 \left(\|p_*\|_2^2+\sum_{l=0}^q\|p_l\|_2^2\right)^{1/2}
 \le e^{C(q+1)^2\log(q+2)}
       (D+1)^{C(q+1)^2}\eta^{-2q}
       \left\|p_*+\sum_{l=0}^q p_lM_l\right\|_2.
\]
Minimal band choices satisfy
$D+1\le C3^q(r+q+2)$.
In particular, if $r,q,\eta^{-1}\le P$ and $P\ge2$, the inverse
norm is at most $P^{C(q+1)^3}$. More generally it is at most
$P^{C(q+1)^3}\eta^{-2q}$ when only $r,q\le P$.
\end{theorem}
\begin{proof}
We track the existing proof, retaining its notation
$w=t(1-t)$, $H_i=\int_0^tP_{d_i}(2u-1)du$,
$\lambda_i=d_i(d_i+1)$, and
$\mathcal L_i C=(wC)''+\lambda_iC$.
The iterated three-band estimate costs at most
\[
 T=C_0^q(D+1)^{6q}\eta^{-q}.
\]
Indeed each iteration is the same absolute-constant three-band lemma;
summing squares does not introduce another dimension-dependent factor.
For $F=p_*+\sum_l p_lM_l$, it bounds every coefficient in
$F=P+\eta\sum_i(H_iA_i+H_i'B_i)$ by $T\|F\|_2$.
The exact primitive formulas give
\[
 -B_i=w\sum_{l=1}^q l p_l\mathcal L_i^{-1}t^{l-1}.
\]

Choose the standard shifted Jacobi polynomials
$J_j(t)=P_j^{(1,1)}(2t-1)$, $0\le j<q$, and put
$\mu_j=(j+1)(j+2)$. Then
$\mathcal L_iJ_j=(\lambda_i-\mu_j)J_j$.
Writing $t^{l-1}=\sum_j a_{jl}J_j$ and
$c_j=\sum_{l=j+1}^q l a_{jl}p_l$ gives
\[
 -B_i=\sum_{j=0}^{q-1}\frac{wJ_jc_j}{\lambda_i-\mu_j}.
\]
All nonzero eigenvalue differences are integers, and
$|\lambda_i-\mu_j|\le3(D+1)^2$.
The Cauchy determinant formula therefore bounds the inverse matrix
by
\[
 \|C^{-1}\|\le q!\,3^{q^2}(D+1)^{2q^2},
\]
up to an additional factor polynomial in $q$ which is harmless below:
the determinant numerator is a nonzero integer, while each entry has
absolute value at most one and each cofactor at most $(q-1)!$.
Consequently all $\|wJ_jc_j\|_2$ are bounded by that quantity times
$qT\|F\|_2$.

Every $J_j$ has its $j$ roots in $(0,1)$ and leading coefficient
$\binom{2j+2}{j}\ge1$. Thus $wJ_j$ satisfies
Lemma~\ref{multiplier} with degree $j+2\le q+1$.
Recovering $c_j$ costs at most
\[
 \sqrt2[4(q+2)(r+1)^2]^{q+2}.
\]
The triangular transformation between $(p_1,\ldots,p_q)$ and the
$c_j$ has inverse norm at most $\exp(Cq^2\log(q+2))$.
For completeness, orthogonality with weight $w$ and
\[
 \int_0^1 wJ_j^2=\frac{j+1}{(2j+3)(j+2)}
\]
give $|a_{jl}|\le C\sqrt{q+1}$. The diagonal coefficient
$a_{l-1,l}$ is the reciprocal of the leading coefficient of $J_{l-1}$,
so it is at least $4^{-l}$. The determinant/cofactor formula, for
entries $l a_{jl}$, gives the asserted inverse bound directly.

We have now bounded all $p_l$, $l\ge1$, by
\[
 e^{C(q+1)^2\log(q+2)}(D+1)^{2q^2+8q+4}\eta^{-q}\|F\|_2,
\]
after increasing $C$ to absorb factors depending only on $q$.
Since $0\le M_l\le1$, subtract those terms. The remainder is
$p_*+p_0M_0$; its trinary representation has polynomial coefficient
$p_*+tp_0$ and coefficients $-p_0$ at each $H_i$.
One further application of the same band estimate costs another $T$.
This proves the stated bound with an absolute $C$.

For minimal bands, writing
$E_i=r+q+1+\sum_{j\le i}(d_j+1)$ gives
$E_i\le3E_{i-1}+C$, hence the displayed bound on $D$.
Finally $\log(D+1)\le Cq+C\log P$ when $r,q\le P$.
All constant factors above can therefore be absorbed into
$P^{C(q+1)^3}$, with an extra $\eta^{-2q}$ if it is not bounded by $P$.
\end{proof}

The same estimates hold for a density $a\le w_0\le b$, with the
explicit additional norm-comparison factors $a^{-1/2},b^{1/2}$.
They also hold on an exact rule for the relevant squared polynomials.
Thus this conditioning step contains no hidden exponential dependence
on the polynomial multiplier degree.
\end{document}
